\section{Preliminaries}
\label{sec:preliminaries}

All target graphs produced by our reductions are finite, simple, and
unweighted. Source graphs may have weights, parallel edges, or loops
when explicitly stated. We write $V(G)$ and $E(G)$ for the vertex and
edge sets of a graph $G$. An \emph{independent set} contains no edge,
including a loop. A \emph{clique} has every pair of distinct vertices adjacent.
A \emph{vertex cover} meets every edge. For an unweighted graph $G$,
let $\alphaI(G)$ denote the maximum independent-set size and
$\vc(G)$ the minimum vertex-cover size.
Taking complements gives
\begin{equation}\label{eq:cover-independent-complement}
 \vc(G)=|V(G)|-\alphaI(G).
\end{equation}
For an optimization instance $I$, write $\OPT(I)$ for its optimum
objective value.
A factor-$r$ approximation for a minimization problem, where $r\ge1$,
returns a feasible solution of cost at most $r\OPT(I)$.
In each reduction the approximation factor and all auxiliary constants
are fixed before the input instance.

For $U\subseteq V(G)$, write $G[U]$ for the subgraph induced by $U$
and $G-U$ for $G[V(G)\setminus U]$. Write $\Delta(G)$ for maximum
degree, $[q]=\{1,\ldots,q\}$ for a positive integer $q$, and
$\binom{U}{2}$ for the set of unordered pairs of distinct elements of $U$.
A spanning subgraph retains all vertices
and some of the edges. A \emph{matching} is a set of edges with
pairwise disjoint endpoints. A matching is \emph{perfect} if it meets
every vertex.

For $S\subseteq V(G)$, the \emph{edge boundary} is
\[
 \cut_G(S)=\{uv\in E(G):u\in S,\ v\in V(G)\setminus S\}.
\]
In a multigraph, this is a multiset, counting parallel edge
occurrences separately.

In a rooted forest $F$, the depth $\operatorname{depth}_F(v)$ of a
vertex $v\in V(F)$ is the number of vertices on its root-to-$v$ path.
Thus roots have depth one. Two vertices are \emph{comparable} if one is
an ancestor of the other. We count a vertex as its own ancestor
and descendant. A \emph{proper} ancestor is a different vertex.
The subtree rooted at a vertex consists of that vertex and all
its descendants. The path from the root of a component to one of its
vertices is an \emph{ancestor path}. We also allow the empty path.
The \emph{restriction} of a forest to a vertex subset deletes all
other vertices and attaches each survivor to its nearest surviving
proper ancestor, or makes it a root if none remains. It preserves
the ancestor order between survivors and never increases their depths.

An \emph{elimination forest} $F$ of a graph $H$ is a rooted forest on
$V(H)$ such that the endpoints of every edge of $H$ are comparable in
the ancestor order. Its height $h(F)$ is the maximum number of vertices on a
root-to-leaf path, with $h(\varnothing)=0$. The \emph{treedepth}
$\td(H)$ is the minimum height of an elimination forest of $H$.
We set $\td(\varnothing)=0$. For
disconnected graphs, treedepth is the maximum treedepth of a component.

The \emph{total depth} of an elimination forest $F$ is
$\totd(F)=\sum_{v\in V(F)}\operatorname{depth}_F(v)$.
Let $\totd(G)$ be the minimum total depth of an
elimination forest of $G$. This objective is also called
\emph{VPT-sum}~\cite{HogemoEtAl2021}. For disconnected graphs,
$\totd$ is the sum of its values on the components. Restricting a forest
to each component proves one inequality, and their disjoint union
proves the other.

A \emph{tree decomposition} of $H$ consists of a tree $T$ and bags
$B_t\subseteq V(H)$ for $t\in V(T)$. For each $v\in V(H)$, write
$T_v=\{t\in V(T):v\in B_t\}$ for the set of nodes whose bags contain
$v$. We require each $T_v$ to be nonempty and to induce a connected
subtree of $T$, and $T_u\cap T_v\ne\varnothing$ for every edge
$uv\in E(H)$. Its width is
$\max_{t\in V(T)}|B_t|-1$, and the \emph{treewidth} $\operatorname{tw}(H)$ is
the minimum such width.

\begin{lemma}[Helly property for subtrees]
\label{lem:subtree-helly}
Let $S_1,\ldots,S_k$ be a nonempty finite family of nonempty
connected subtrees of a tree $T$. If $S_i\cap S_j\ne\varnothing$
for every $i,j\in[k]$, then $\bigcap_{i=1}^k S_i\ne\varnothing$.
\end{lemma}

\begin{proof}
Root $T$, and let $r_i$ be the vertex of $S_i$ nearest the root.
Connectedness makes $r_i$ unique and an ancestor of every vertex
of $S_i$. Choose $r_j$ of maximum depth and put $z=r_j$.
For each $i\in[k]$, choose $x\in S_i\cap S_j$. Both $r_i$ and $z$ are
ancestors of $x$, and $z$ is at least as deep as $r_i$. Thus $z$
lies on the $r_i$--$x$ path, which is contained in $S_i$.
Hence $z$ belongs to every $S_i$.
\end{proof}

\begin{lemma}[Three paths in a tree]
\label{lem:three-tree-paths}
For any three vertices $a,b,c$ of a tree, not necessarily distinct,
the three paths joining their pairs have exactly one common vertex.
Here the path from a vertex to itself consists of that vertex.
\end{lemma}

\begin{proof}
Let $z$ be the first vertex on the $c$--$a$ path that belongs to
the $a$--$b$ path. Uniqueness of paths forces the $c$--$b$ path
to consist of the $c$--$z$ path followed by the $z$--$b$ path.
The $z$--$a$ and $z$--$b$ portions of the $a$--$b$ path have only
$z$ in common, so the three paths intersect exactly at $z$.
\end{proof}

\begin{lemma}[Minor monotonicity of treedepth and treewidth]
\label{lem:minor-monotonicity}
If $H'$ is obtained from $H$ by vertex deletions, edge deletions,
and edge contractions, then
$\td(H')\leq\td(H)$ and
$\operatorname{tw}(H')\leq\operatorname{tw}(H)$.
\end{lemma}

\begin{proof}
For treedepth, vertex deletion is handled by restricting an
elimination forest, and edge deletion leaves it valid. For an edge
contraction $uv$, its endpoints are comparable in any elimination
forest. Suppose $u$ is above $v$. Restrict the forest to
$V(H)\setminus\{v\}$, using $u$ for the contracted vertex. All
edges not incident with $v$ retain comparable endpoints. Every
ancestor of $v$ is comparable with $u$, and every descendant of $v$
is a descendant of $u$, so the edges transferred from $v$ to $u$
also have comparable endpoints. Restriction does not increase height.

For treewidth, delete vertices from every bag or leave the bags
unchanged when deleting edges. To contract $uv$, replace $u$ and $v$
by the contracted vertex in every bag. The node sets $T_u,T_v$
intersect because $uv$ is an edge, so their union is connected. All
edges remain covered and no bag grows.
\end{proof}

\begin{lemma}[An ancestor in a connected set]
\label{lem:connected-set-ancestor}
Let $F$ be an elimination forest of $H$, and let
$\varnothing\ne U\subseteq V(H)$ induce a connected subgraph.
The vertex of $U$ with minimum depth in $F$ is unique and is an
ancestor of every vertex of $U$.
\end{lemma}

\begin{proof}
Choose a minimum-depth vertex $r\in U$. If a path in $H[U]$ starting
at $r$ left its forest subtree, the first edge leaving it would join
a descendant of $r$ to a proper ancestor of $r$, since its endpoints
are comparable. This contradicts the choice of $r$. Thus all of $U$
lies below $r$, which also proves uniqueness.
\end{proof}
